Digital signatures underpin firmware updates, secure boot, messaging and
public-key infrastructures, and the schemes deployed today rest on integer
factorisation or on the Elliptic Curve Discrete Logarithm Problem, both of
which a large quantum computer would solve with Shor's algorithm. Among the
most promising replacements are signatures obtained by applying the
Fiat--Shamir transform~\cite{DagdelenFG13,FiatS86} to a zero-knowledge
identification protocol repeated in parallel --- LESS~\cite{LESS-spec},
CROSS~\cite{CROSSspec} and MEDS~\cite{ChouNPRRST23,MEDS_PQC} among them. Their
weak point is signature size, and the standard remedy is
\GGM-tree seed compression: the $t$ ephemeral secrets are derived from a single
root seed, so that the signer publishes a handful of internal nodes rather than
$t$ independent seeds.

During signing, only the seeds of the revealed rounds may be released; the
challenge-dependent ones must stay hidden. This is the single point at which
zero knowledge is enforced in the implementation, and recent work has shown it
to be a fault-attack surface: a fault in the seed-publication logic can leak a
hidden seed alongside the zero-knowledge response of the same round, which is
enough for key recovery or universal
forgery~\cite{FaultToForge2025,ZKFault}.


%In 2022, the National Institute 
%of Standards and Technology selected the first set of post-quantum 
%digital signature schemes as part of its standardization process, 
%marking a major step toward the deployment of quantum-resistant 
%cryptography.





\paragraph{Why seed trees, and not one scheme.}
The three schemes named above have all now left the NIST additional-signatures
process: MEDS was not selected for the second round, and CROSS and LESS were
eliminated at the end of it~\cite{NISTIR8610}. It would be a mistake to read
this as closing the subject. The compression technique is not specific to
those designs: it is the standard way to shrink any Fiat--Shamir signature
built from a parallel-repeated identification protocol, and FAEST, SDitH and
MQOM --- all of which did advance --- publish an all-but-one vector commitment
derived from the same tree. What we analyse is the mechanism, not its host. We
attack MEDS because its reference implementation realises that mechanism in
its cleanest form, and the conclusions transfer to any scheme publishing a
seed-tree path. The timing is favourable rather than late: the third round is
explicitly an implementation-security phase.

\subsection{Related Work}
\label{sec:relwork}
Fault attacks on seed-tree-based post-quantum signatures are an active line of
work. ZKFault~\cite{ZKFault} faults the seed publication of LESS-v1, extended
to CROSS, disclosing hidden ephemeral seeds and recovering secret material.
Fault-to-Forge~\cite{FaultToForge2025} revisits LESS-v2, whose incomplete tree
defeats that strategy, and finds new injection points in tree generation and
seed publication, showing that secret-equivalent information already suffices
for universal forgery. The attack of~\cite{HuangHHYSTY26} avoids the tree
altogether, using iteration-skip and loop-abort faults in the
digest-preprocessing step that precedes tree generation. Elsewhere in the
family,~\cite{cryptoeprint:2026/1266} faults tree generation in
PERK~\cite{BettaiebBDEGKP24},~\cite{JendralD25} targets FAEST,
and~\cite{BandaBK26} targets Mirath and RYDE.

We unify these results under a single Generic ZK Seed Tree abstraction,
formalising the invariants they implicitly exploit, and add practical fault
attacks against MEDS.

\paragraph{Our Contributions.}
This work provides a unified analysis of fault attacks against
post-quantum signature schemes relying on seed-tree compression.
Our main contributions are as follows:
\begin{itemize}

\item \textbf{Generic Seed-Tree Abstraction.}
We introduce the \emph{Generic ZK Seed Tree} (GZKST), a unified model
capturing the seed-tree generation mechanisms used in schemes such as
LESS-v2, MEDS, and CROSS.

\item \textbf{Unified View of Existing Attacks.}
We show that previously proposed fault attacks against LESS-v1 and
LESS-v2 can be interpreted as violations of a common seed-hiding
invariant.

\item \textbf{A Taxonomy of Fault Surfaces.}
We classify every published fault attack on this family by which component
of the abstraction it corrupts (Section~\ref{sec:taxonomy}), showing that
five independent attacks against four schemes occupy only three of six
components. The empty cells are predictions, and one of them led us
directly to a new attack surface in the MEDS reference implementation
that has no analogue in prior work.

\item \textbf{Query Complexity Analysis.}
We derive bounds on the number of faulted signatures required for
secret-key recovery and show that only a small number of successful
queries is needed for all MEDS parameter sets.

\item \textbf{Practical Attacks on MEDS.}
We show fault attacks against the MEDS reference
implementation on an ARM Cortex-M4 using a ChipWhisperer-Lite,
revealing several fault surfaces that enable complete or partial
seed-tree disclosure.

\item \textbf{Countermeasures for MEDS, and their limits.}
We propose and evaluate three families of countermeasures. We also show
that the obvious defence --- verifying a signature before releasing it ---
is weaker than it looks against this particular implementation, because
path construction and path reconstruction share a single function, so the
same glitch at the same code offset defeats both passes
(Section~\ref{sec:counter}).

\end{itemize}

